% Generated draft: requires mathematical and source review.
\documentclass[11pt]{article}
\usepackage[margin=25mm]{geometry}
\usepackage{fontspec}
\setmainfont{Latin Modern Roman}
\usepackage{amsmath,amssymb,mathtools,unicode-math}
\setmathfont{Latin Modern Math}
\usepackage{enumitem,needspace,xurl,hyperref}
\setlength{\parindent}{0pt}
\setlength{\parskip}{5pt}
\setlength{\emergencystretch}{3em}
\newcommand{\N}{\mathbb{N}}
\newcommand{\Z}{\mathbb{Z}}
\newcommand{\Q}{\mathbb{Q}}
\newcommand{\R}{\mathbb{R}}
\newcommand{\C}{\mathbb{C}}
\newcommand{\F}{\mathbb{F}}
\newcommand{\A}{\mathbb{A}}
\newcommand{\PP}{\mathbb{P}}
\newcommand{\E}{\mathbb{E}}
\newcommand{\eps}{\varepsilon}
\newcommand{\dd}{\,\mathrm{d}}
\begin{document}
\section*{UnsolvedMath ID: 30000977}
\textbf{Source problem code:} catalogue entry
\subsection{Definitions and mathematical statement}
In the usual characteristic-zero projective-surface setting, let \(X\) be a normal projective surface of general type with canonical singularities, and assume that its canonical divisor \(K_X\) is ample and base-point-free. Write \(p_g(X)=h^0(X,\omega_X)\). The complete canonical linear system defines a morphism
\[
 \varphi=\varphi_{|K_X|}:X\longrightarrow
 W\subseteq\PP^{p_g(X)-1}.
\]
The source assumes that the image \(W\) is a singular surface of minimal degree. Minimal degree means, in its nondegenerate projective embedding,
\[
 \deg W=\operatorname{codim}_{\PP^{p_g(X)-1}}W+1=p_g(X)-2.
\]
The condition singular is essential; the question is not restricted to a smooth minimal-degree image or to one chosen singular cone.

Since \(\varphi^*\mathcal O_W(1)=\omega_X\) is ample, a curve cannot be contracted by \(\varphi\); thus the canonical map is finite onto the given surface image in this setting. Its degree is the generic degree
\[
 \deg\varphi=[k(X):k(W)],
\]
not the projective degree of \(W\). The extracted question is the universal bound
\[
 \forall\,(X,\varphi,W)\text{ satisfying these hypotheses},\qquad
 \deg\varphi\leq4\ ?
\]
Equivalently, under the stated finite-map intersection convention,
\(K_X^2=(\deg\varphi)\deg W=(\deg\varphi)(p_g(X)-2)\), so the same requested bound is \(K_X^2\leq4(p_g(X)-2)\). No Galois, cyclic, flatness, regularity, or smoothness assumption on the cover has been added.

The ground field and characteristic are not explicit in the record; the characteristic-zero setting above is identified as the usual interpretation, not an additional verified report hypothesis. The supplied assessment dated 2026-08-21 describes the extraction as a conditional fragment referring to a preceding canonical-cover question, whereas the supplied original and cleaned text explicitly state the degree-bound question with the displayed hypotheses. That discrepancy must be checked for omitted context. The assessment reports the underlying bound as a research question, not a theorem.
\subsection{Short English statement}
Must a finite canonical map from a canonically singular general-type surface with ample base-point-free canonical divisor have degree at most four when its image is a singular surface of minimal degree? The bound concerns the map, and no Galois-cover restriction is supplied.
\subsection{Sources}
\begin{itemize}
\item Local UnsolvedMath registry, permanent ID \(30000977\), exact source code \texttt{catalogue entry}, proposed year 2008; no individual proposer is supplied.
\item \emph{Classical Algebraic Geometry (2009)}, Oberwolfach Reports, DOI 10.4171/owr/2008/26, pp. 1425--1474. \url{https://doi.org/10.4171/owr/2008/26}.
\item Supplied reference labelled ``recent canonical-cover literature'': \url{https://link.springer.com/article/10.1007/s13163-023-00462-5}.
\item The supplied \texttt{corrected\_verified} note reports editorial and LaTeX cleanup without added source-dependent content. The open label and assessment checked 2026-08-21 are dated supplied triage; its conditional-fragment description and the scope of the recent reference require independent source review.
\end{itemize}
\end{document}
